\section{Kết quả đầy đủ của benchmark nhận diện trên điện thoại}\label{sec:phone-full-results}
Phần này công bố đủ 32 cấu hình của bộ benchmark có provenance, làm căn cứ cho ba bảng tóm tắt và kết luận tại Mục~\ref{sec:phone-benchmark-results}. Không loại cấu hình kém thuận lợi. Mỗi cấu hình dùng cùng 256 ảnh test, gồm 86 non\_flood, 56 low, 60 medium và 54 high; mỗi ảnh có một dự đoán trong từng phiên. Các phiên dùng lại ảnh nên không làm tăng số ảnh độc lập lên 8.192.

Mã cấu hình có dạng \texttt{seed-biến thể-máy-runtime}. B là baseline, P là cắt tỉa; RM15 là realme 15, RN11 là Redmi Note 11, RN13 là Redmi Note 13 và S10e là Samsung S10e. O là ONNX, T là ExecuTorch/PTE. Ví dụ, \texttt{1024-P-RN13-T} là seed 1024 cắt tỉa chạy PTE trên Redmi Note 13. Thông tin thiết bị, dataset và phạm vi đo được mô tả ở Mục~\ref{sec:phone-benchmark-setup}.

\subsection{Chất lượng phân loại theo cấu hình}
Accuracy (Acc.) là tỷ lệ dự đoán đúng; P/R/F1 macro là trung bình có trọng số bằng nhau của chỉ số tương ứng ở bốn lớp. Macro-Recall cũng chính là balanced accuracy trong bài toán này. Các giá trị phần trăm được làm tròn hai chữ số; CSV giữ độ chính xác cao hơn.
% Generated by scripts/build_phone_report_tables.py; do not edit.
\begingroup\small\setlength{\tabcolsep}{4pt}
\begin{longtable}{lrrrr}
\caption{Chất lượng phân loại của đủ 32 cấu hình điện thoại (\%)}\label{tab:phone-full-quality}\\
\toprule Cấu hình & Acc. & P macro & R macro & F1 macro \\ \midrule\endfirsthead
\multicolumn{5}{l}{\tablename~\thetable{} (tiếp)}\\
\toprule Cấu hình & Acc. & P macro & R macro & F1 macro \\ \midrule\endhead
\midrule\multicolumn{5}{r}{Còn tiếp ở trang sau}\\\endfoot
\bottomrule\endlastfoot
159-B-RM15-O & 75,39 & 73,98 & 73,39 & 73,60 \\
159-B-RM15-T & 75,39 & 73,98 & 73,39 & 73,60 \\
159-B-RN11-O & 74,61 & 73,24 & 72,48 & 72,73 \\
159-B-RN11-T & 74,61 & 73,24 & 72,48 & 72,73 \\
159-B-RN13-O & 75,00 & 73,64 & 72,93 & 73,18 \\
159-B-RN13-T & 75,00 & 73,64 & 72,93 & 73,18 \\
159-B-S10e-O & 75,39 & 74,20 & 73,35 & 73,65 \\
159-B-S10e-T & 75,39 & 74,20 & 73,35 & 73,65 \\
159-P-RM15-O & 76,95 & 74,60 & 75,13 & 74,58 \\
159-P-RM15-T & 76,95 & 74,60 & 75,13 & 74,58 \\
159-P-RN11-O & 76,56 & 74,24 & 74,84 & 74,36 \\
159-P-RN11-T & 76,56 & 74,24 & 74,84 & 74,36 \\
159-P-RN13-O & 76,95 & 74,57 & 75,13 & 74,62 \\
159-P-RN13-T & 76,95 & 74,57 & 75,13 & 74,62 \\
159-P-S10e-O & 75,78 & 73,35 & 73,93 & 73,49 \\
159-P-S10e-T & 75,78 & 73,35 & 73,93 & 73,49 \\
1024-B-RM15-O & 75,39 & 74,37 & 73,50 & 73,83 \\
1024-B-RM15-T & 75,39 & 74,37 & 73,50 & 73,83 \\
1024-B-RN11-O & 75,39 & 74,33 & 73,50 & 73,82 \\
1024-B-RN11-T & 75,39 & 74,33 & 73,50 & 73,82 \\
1024-B-RN13-O & 75,78 & 74,68 & 73,96 & 74,25 \\
1024-B-RN13-T & 75,78 & 74,68 & 73,96 & 74,25 \\
1024-B-S10e-O & 76,17 & 75,27 & 74,38 & 74,72 \\
1024-B-S10e-T & 76,17 & 75,27 & 74,38 & 74,72 \\
1024-P-RM15-O & 76,95 & 74,79 & 75,21 & 74,89 \\
1024-P-RM15-T & 76,95 & 74,79 & 75,21 & 74,89 \\
1024-P-RN11-O & 76,17 & 74,04 & 74,28 & 74,10 \\
1024-P-RN11-T & 76,17 & 74,04 & 74,28 & 74,10 \\
1024-P-RN13-O & 76,56 & 74,30 & 74,79 & 74,40 \\
1024-P-RN13-T & 76,56 & 74,30 & 74,79 & 74,40 \\
1024-P-S10e-O & 76,95 & 74,86 & 75,21 & 74,90 \\
1024-P-S10e-T & 76,95 & 74,86 & 75,21 & 74,90 \\
\end{longtable}\endgroup


\subsection{Sai lệch thứ bậc và lỗi nặng}
Đặt $s(\mathrm{non\_flood})=0$, $s(\mathrm{low})=1$, $s(\mathrm{medium})=2$, $s(\mathrm{high})=3$. Với $e_i=s(\hat y_i)-s(y_i)$ và $n=256$, hai chỉ số bổ sung được tính như sau:
\[
\operatorname{MAE}=\frac{1}{n}\sum_{i=1}^{n}|e_i|,
\qquad
\operatorname{RMSE}=\sqrt{\frac{1}{n}\sum_{i=1}^{n}e_i^2}.
\]
Lỗi nặng là số ảnh có $|e_i|\geq2$. Khoảng cách giữa các mức được quy ước bằng một, không phải độ sâu ngập hay chi phí cứu hộ. MAE/RMSE chưa có trong JSON gốc; chúng được tính lại từ dự đoán, không chạy suy luận mới và không sửa dữ liệu nguồn.
% Generated by scripts/build_phone_report_tables.py; do not edit.
\begingroup\small\setlength{\tabcolsep}{4pt}
\begin{longtable}{lrrr}
\caption{Sai lệch thứ bậc của đủ 32 cấu hình điện thoại}\label{tab:phone-full-ordinal}\\
\toprule Cấu hình & MAE (mức) & RMSE (mức) & Lỗi nặng/256 \\ \midrule\endfirsthead
\multicolumn{4}{l}{\tablename~\thetable{} (tiếp)}\\
\toprule Cấu hình & MAE (mức) & RMSE (mức) & Lỗi nặng/256 \\ \midrule\endhead
\midrule\multicolumn{4}{r}{Còn tiếp ở trang sau}\\\endfoot
\bottomrule\endlastfoot
159-B-RM15-O & 0,2656 & 0,5590 & 4 \\
159-B-RM15-T & 0,2656 & 0,5590 & 4 \\
159-B-RN11-O & 0,2734 & 0,5660 & 4 \\
159-B-RN11-T & 0,2734 & 0,5660 & 4 \\
159-B-RN13-O & 0,2695 & 0,5625 & 4 \\
159-B-RN13-T & 0,2695 & 0,5625 & 4 \\
159-B-S10e-O & 0,2656 & 0,5590 & 4 \\
159-B-S10e-T & 0,2656 & 0,5590 & 4 \\
159-P-RM15-O & 0,2500 & 0,5449 & 4 \\
159-P-RM15-T & 0,2500 & 0,5449 & 4 \\
159-P-RN11-O & 0,2461 & 0,5266 & 2 \\
159-P-RN11-T & 0,2461 & 0,5266 & 2 \\
159-P-RN13-O & 0,2461 & 0,5340 & 3 \\
159-P-RN13-T & 0,2461 & 0,5340 & 3 \\
159-P-S10e-O & 0,2578 & 0,5449 & 3 \\
159-P-S10e-T & 0,2578 & 0,5449 & 3 \\
1024-B-RM15-O & 0,2539 & 0,5192 & 2 \\
1024-B-RM15-T & 0,2539 & 0,5192 & 2 \\
1024-B-RN11-O & 0,2578 & 0,5303 & 3 \\
1024-B-RN11-T & 0,2578 & 0,5303 & 3 \\
1024-B-RN13-O & 0,2500 & 0,5154 & 2 \\
1024-B-RN13-T & 0,2500 & 0,5154 & 2 \\
1024-B-S10e-O & 0,2461 & 0,5116 & 2 \\
1024-B-S10e-T & 0,2461 & 0,5116 & 2 \\
1024-P-RM15-O & 0,2422 & 0,5154 & 3 \\
1024-P-RM15-T & 0,2422 & 0,5154 & 3 \\
1024-P-RN11-O & 0,2539 & 0,5340 & 4 \\
1024-P-RN11-T & 0,2539 & 0,5340 & 4 \\
1024-P-RN13-O & 0,2461 & 0,5192 & 3 \\
1024-P-RN13-T & 0,2461 & 0,5192 & 3 \\
1024-P-S10e-O & 0,2422 & 0,5154 & 3 \\
1024-P-S10e-T & 0,2422 & 0,5154 & 3 \\
\end{longtable}\endgroup


\subsection{Thời gian xử lý theo cấu hình}
Mean là trung bình thời gian từng ảnh; P50/P95 dùng nearest-rank sau khi loại năm lượt warmup khỏi latency. Đơn vị là ms; phạm vi là giải mã, tiền xử lý và suy luận trong \texttt{classifyImage}. Bảng giữ thời gian riêng từng máy, không dùng thời gian tổng hợp thay cho kết quả từng thiết bị.
% Generated by scripts/build_phone_report_tables.py; do not edit.
\begingroup\small\setlength{\tabcolsep}{4pt}
\begin{longtable}{lrrr}
\caption{Thời gian xử lý ảnh của đủ 32 cấu hình điện thoại (ms)}\label{tab:phone-full-latency}\\
\toprule Cấu hình & Mean & P50 & P95 \\ \midrule\endfirsthead
\multicolumn{4}{l}{\tablename~\thetable{} (tiếp)}\\
\toprule Cấu hình & Mean & P50 & P95 \\ \midrule\endhead
\midrule\multicolumn{4}{r}{Còn tiếp ở trang sau}\\\endfoot
\bottomrule\endlastfoot
159-B-RM15-O & 85,218 & 72,280 & 155,767 \\
159-B-RM15-T & 45,521 & 37,836 & 97,626 \\
159-B-RN11-O & 178,741 & 177,417 & 253,728 \\
159-B-RN11-T & 115,493 & 110,716 & 206,998 \\
159-B-RN13-O & 202,633 & 188,961 & 296,778 \\
159-B-RN13-T & 87,676 & 73,750 & 176,585 \\
159-B-S10e-O & 84,328 & 71,242 & 167,766 \\
159-B-S10e-T & 41,277 & 32,301 & 94,351 \\
159-P-RM15-O & 80,660 & 73,644 & 137,510 \\
159-P-RM15-T & 49,079 & 41,027 & 103,708 \\
159-P-RN11-O & 178,418 & 172,520 & 278,802 \\
159-P-RN11-T & 106,830 & 100,001 & 187,512 \\
159-P-RN13-O & 193,530 & 176,095 & 289,758 \\
159-P-RN13-T & 90,688 & 76,153 & 184,379 \\
159-P-S10e-O & 84,116 & 71,372 & 165,529 \\
159-P-S10e-T & 48,430 & 38,978 & 115,327 \\
1024-B-RM15-O & 82,673 & 74,863 & 145,783 \\
1024-B-RM15-T & 45,693 & 38,414 & 96,905 \\
1024-B-RN11-O & 179,317 & 169,647 & 275,742 \\
1024-B-RN11-T & 102,287 & 92,059 & 201,230 \\
1024-B-RN13-O & 190,192 & 175,082 & 285,861 \\
1024-B-RN13-T & 87,984 & 73,436 & 178,710 \\
1024-B-S10e-O & 84,654 & 72,211 & 171,631 \\
1024-B-S10e-T & 40,603 & 31,508 & 95,932 \\
1024-P-RM15-O & 81,943 & 74,471 & 135,131 \\
1024-P-RM15-T & 48,944 & 39,344 & 103,338 \\
1024-P-RN11-O & 177,966 & 172,819 & 274,626 \\
1024-P-RN11-T & 102,766 & 90,892 & 202,812 \\
1024-P-RN13-O & 246,763 & 239,020 & 335,798 \\
1024-P-RN13-T & 88,475 & 73,147 & 180,651 \\
1024-P-S10e-O & 86,190 & 73,526 & 163,776 \\
1024-P-S10e-T & 43,270 & 34,401 & 103,429 \\
\end{longtable}\endgroup


\subsection{Precision, Recall và F1 theo từng lớp}
Bảng gồm 128 hàng, tương ứng 32 cấu hình nhân bốn lớp. Cột số ảnh là số mẫu thật của lớp trong test; không phải số ảnh model dự đoán vào lớp đó. F1 được tính riêng từ Precision/Recall từng lớp, sau đó mới lấy trung bình để có macro-F1. Việc hai runtime có cùng nhãn top-1 làm các chỉ số này bằng nhau trong từng cặp máy/model, nhưng không chứng minh vector xác suất của hai runtime giống tuyệt đối.
% Generated by scripts/build_phone_report_tables.py; do not edit.
\begingroup\small\setlength{\tabcolsep}{4pt}
\begin{longtable}{lrrrrr}
\caption{Chỉ số từng lớp của đủ 32 cấu hình điện thoại (\%)}\label{tab:phone-full-classes}\\
\toprule Cấu hình & Lớp & Số ảnh & Precision & Recall & F1 \\ \midrule\endfirsthead
\multicolumn{6}{l}{\tablename~\thetable{} (tiếp)}\\
\toprule Cấu hình & Lớp & Số ảnh & Precision & Recall & F1 \\ \midrule\endhead
\midrule\multicolumn{6}{r}{Còn tiếp ở trang sau}\\\endfoot
\bottomrule\endlastfoot
159-B-RM15-O & non\_flood & 86 & 84,95 & 91,86 & 88,27 \\
159-B-RM15-O & low & 56 & 64,71 & 58,93 & 61,68 \\
159-B-RM15-O & medium & 60 & 63,93 & 65,00 & 64,46 \\
159-B-RM15-O & high & 54 & 82,35 & 77,78 & 80,00 \\
159-B-RM15-T & non\_flood & 86 & 84,95 & 91,86 & 88,27 \\
159-B-RM15-T & low & 56 & 64,71 & 58,93 & 61,68 \\
159-B-RM15-T & medium & 60 & 63,93 & 65,00 & 64,46 \\
159-B-RM15-T & high & 54 & 82,35 & 77,78 & 80,00 \\
159-B-RN11-O & non\_flood & 86 & 84,04 & 91,86 & 87,78 \\
159-B-RN11-O & low & 56 & 64,00 & 57,14 & 60,38 \\
159-B-RN11-O & medium & 60 & 62,90 & 65,00 & 63,93 \\
159-B-RN11-O & high & 54 & 82,00 & 75,93 & 78,85 \\
159-B-RN11-T & non\_flood & 86 & 84,04 & 91,86 & 87,78 \\
159-B-RN11-T & low & 56 & 64,00 & 57,14 & 60,38 \\
159-B-RN11-T & medium & 60 & 62,90 & 65,00 & 63,93 \\
159-B-RN11-T & high & 54 & 82,00 & 75,93 & 78,85 \\
159-B-RN13-O & non\_flood & 86 & 84,95 & 91,86 & 88,27 \\
159-B-RN13-O & low & 56 & 64,71 & 58,93 & 61,68 \\
159-B-RN13-O & medium & 60 & 62,90 & 65,00 & 63,93 \\
159-B-RN13-O & high & 54 & 82,00 & 75,93 & 78,85 \\
159-B-RN13-T & non\_flood & 86 & 84,95 & 91,86 & 88,27 \\
159-B-RN13-T & low & 56 & 64,71 & 58,93 & 61,68 \\
159-B-RN13-T & medium & 60 & 62,90 & 65,00 & 63,93 \\
159-B-RN13-T & high & 54 & 82,00 & 75,93 & 78,85 \\
159-B-S10e-O & non\_flood & 86 & 84,95 & 91,86 & 88,27 \\
159-B-S10e-O & low & 56 & 64,71 & 58,93 & 61,68 \\
159-B-S10e-O & medium & 60 & 63,49 & 66,67 & 65,04 \\
159-B-S10e-O & high & 54 & 83,67 & 75,93 & 79,61 \\
159-B-S10e-T & non\_flood & 86 & 84,95 & 91,86 & 88,27 \\
159-B-S10e-T & low & 56 & 64,71 & 58,93 & 61,68 \\
159-B-S10e-T & medium & 60 & 63,49 & 66,67 & 65,04 \\
159-B-S10e-T & high & 54 & 83,67 & 75,93 & 79,61 \\
159-P-RM15-O & non\_flood & 86 & 91,01 & 94,19 & 92,57 \\
159-P-RM15-O & low & 56 & 65,45 & 64,29 & 64,86 \\
159-P-RM15-O & medium & 60 & 67,35 & 55,00 & 60,55 \\
159-P-RM15-O & high & 54 & 74,60 & 87,04 & 80,34 \\
159-P-RM15-T & non\_flood & 86 & 91,01 & 94,19 & 92,57 \\
159-P-RM15-T & low & 56 & 65,45 & 64,29 & 64,86 \\
159-P-RM15-T & medium & 60 & 67,35 & 55,00 & 60,55 \\
159-P-RM15-T & high & 54 & 74,60 & 87,04 & 80,34 \\
159-P-RN11-O & non\_flood & 86 & 90,91 & 93,02 & 91,95 \\
159-P-RN11-O & low & 56 & 64,29 & 64,29 & 64,29 \\
159-P-RN11-O & medium & 60 & 64,71 & 55,00 & 59,46 \\
159-P-RN11-O & high & 54 & 77,05 & 87,04 & 81,74 \\
159-P-RN11-T & non\_flood & 86 & 90,91 & 93,02 & 91,95 \\
159-P-RN11-T & low & 56 & 64,29 & 64,29 & 64,29 \\
159-P-RN11-T & medium & 60 & 64,71 & 55,00 & 59,46 \\
159-P-RN11-T & high & 54 & 77,05 & 87,04 & 81,74 \\
159-P-RN13-O & non\_flood & 86 & 91,01 & 94,19 & 92,57 \\
159-P-RN13-O & low & 56 & 65,45 & 64,29 & 64,86 \\
159-P-RN13-O & medium & 60 & 66,00 & 55,00 & 60,00 \\
159-P-RN13-O & high & 54 & 75,81 & 87,04 & 81,03 \\
159-P-RN13-T & non\_flood & 86 & 91,01 & 94,19 & 92,57 \\
159-P-RN13-T & low & 56 & 65,45 & 64,29 & 64,86 \\
159-P-RN13-T & medium & 60 & 66,00 & 55,00 & 60,00 \\
159-P-RN13-T & high & 54 & 75,81 & 87,04 & 81,03 \\
159-P-S10e-O & non\_flood & 86 & 90,91 & 93,02 & 91,95 \\
159-P-S10e-O & low & 56 & 63,64 & 62,50 & 63,06 \\
159-P-S10e-O & medium & 60 & 63,46 & 55,00 & 58,93 \\
159-P-S10e-O & high & 54 & 75,41 & 85,19 & 80,00 \\
159-P-S10e-T & non\_flood & 86 & 90,91 & 93,02 & 91,95 \\
159-P-S10e-T & low & 56 & 63,64 & 62,50 & 63,06 \\
159-P-S10e-T & medium & 60 & 63,46 & 55,00 & 58,93 \\
159-P-S10e-T & high & 54 & 75,41 & 85,19 & 80,00 \\
1024-B-RM15-O & non\_flood & 86 & 89,66 & 90,70 & 90,17 \\
1024-B-RM15-O & low & 56 & 61,82 & 60,71 & 61,26 \\
1024-B-RM15-O & medium & 60 & 60,61 & 66,67 & 63,49 \\
1024-B-RM15-O & high & 54 & 85,42 & 75,93 & 80,39 \\
1024-B-RM15-T & non\_flood & 86 & 89,66 & 90,70 & 90,17 \\
1024-B-RM15-T & low & 56 & 61,82 & 60,71 & 61,26 \\
1024-B-RM15-T & medium & 60 & 60,61 & 66,67 & 63,49 \\
1024-B-RM15-T & high & 54 & 85,42 & 75,93 & 80,39 \\
1024-B-RN11-O & non\_flood & 86 & 89,66 & 90,70 & 90,17 \\
1024-B-RN11-O & low & 56 & 60,71 & 60,71 & 60,71 \\
1024-B-RN11-O & medium & 60 & 61,54 & 66,67 & 64,00 \\
1024-B-RN11-O & high & 54 & 85,42 & 75,93 & 80,39 \\
1024-B-RN11-T & non\_flood & 86 & 89,66 & 90,70 & 90,17 \\
1024-B-RN11-T & low & 56 & 60,71 & 60,71 & 60,71 \\
1024-B-RN11-T & medium & 60 & 61,54 & 66,67 & 64,00 \\
1024-B-RN11-T & high & 54 & 85,42 & 75,93 & 80,39 \\
1024-B-RN13-O & non\_flood & 86 & 89,66 & 90,70 & 90,17 \\
1024-B-RN13-O & low & 56 & 61,82 & 60,71 & 61,26 \\
1024-B-RN13-O & medium & 60 & 61,54 & 66,67 & 64,00 \\
1024-B-RN13-O & high & 54 & 85,71 & 77,78 & 81,55 \\
1024-B-RN13-T & non\_flood & 86 & 89,66 & 90,70 & 90,17 \\
1024-B-RN13-T & low & 56 & 61,82 & 60,71 & 61,26 \\
1024-B-RN13-T & medium & 60 & 61,54 & 66,67 & 64,00 \\
1024-B-RN13-T & high & 54 & 85,71 & 77,78 & 81,55 \\
1024-B-S10e-O & non\_flood & 86 & 89,66 & 90,70 & 90,17 \\
1024-B-S10e-O & low & 56 & 61,82 & 60,71 & 61,26 \\
1024-B-S10e-O & medium & 60 & 62,12 & 68,33 & 65,08 \\
1024-B-S10e-O & high & 54 & 87,50 & 77,78 & 82,35 \\
1024-B-S10e-T & non\_flood & 86 & 89,66 & 90,70 & 90,17 \\
1024-B-S10e-T & low & 56 & 61,82 & 60,71 & 61,26 \\
1024-B-S10e-T & medium & 60 & 62,12 & 68,33 & 65,08 \\
1024-B-S10e-T & high & 54 & 87,50 & 77,78 & 82,35 \\
1024-P-RM15-O & non\_flood & 86 & 88,89 & 93,02 & 90,91 \\
1024-P-RM15-O & low & 56 & 62,26 & 58,93 & 60,55 \\
1024-P-RM15-O & medium & 60 & 66,67 & 60,00 & 63,16 \\
1024-P-RM15-O & high & 54 & 81,36 & 88,89 & 84,96 \\
1024-P-RM15-T & non\_flood & 86 & 88,89 & 93,02 & 90,91 \\
1024-P-RM15-T & low & 56 & 62,26 & 58,93 & 60,55 \\
1024-P-RM15-T & medium & 60 & 66,67 & 60,00 & 63,16 \\
1024-P-RM15-T & high & 54 & 81,36 & 88,89 & 84,96 \\
1024-P-RN11-O & non\_flood & 86 & 88,89 & 93,02 & 90,91 \\
1024-P-RN11-O & low & 56 & 61,11 & 58,93 & 60,00 \\
1024-P-RN11-O & medium & 60 & 65,45 & 60,00 & 62,61 \\
1024-P-RN11-O & high & 54 & 80,70 & 85,19 & 82,88 \\
1024-P-RN11-T & non\_flood & 86 & 88,89 & 93,02 & 90,91 \\
1024-P-RN11-T & low & 56 & 61,11 & 58,93 & 60,00 \\
1024-P-RN11-T & medium & 60 & 65,45 & 60,00 & 62,61 \\
1024-P-RN11-T & high & 54 & 80,70 & 85,19 & 82,88 \\
1024-P-RN13-O & non\_flood & 86 & 88,89 & 93,02 & 90,91 \\
1024-P-RN13-O & low & 56 & 62,26 & 58,93 & 60,55 \\
1024-P-RN13-O & medium & 60 & 66,04 & 58,33 & 61,95 \\
1024-P-RN13-O & high & 54 & 80,00 & 88,89 & 84,21 \\
1024-P-RN13-T & non\_flood & 86 & 88,89 & 93,02 & 90,91 \\
1024-P-RN13-T & low & 56 & 62,26 & 58,93 & 60,55 \\
1024-P-RN13-T & medium & 60 & 66,04 & 58,33 & 61,95 \\
1024-P-RN13-T & high & 54 & 80,00 & 88,89 & 84,21 \\
1024-P-S10e-O & non\_flood & 86 & 87,91 & 93,02 & 90,40 \\
1024-P-S10e-O & low & 56 & 62,26 & 58,93 & 60,55 \\
1024-P-S10e-O & medium & 60 & 67,92 & 60,00 & 63,72 \\
1024-P-S10e-O & high & 54 & 81,36 & 88,89 & 84,96 \\
1024-P-S10e-T & non\_flood & 86 & 87,91 & 93,02 & 90,40 \\
1024-P-S10e-T & low & 56 & 62,26 & 58,93 & 60,55 \\
1024-P-S10e-T & medium & 60 & 67,92 & 60,00 & 63,72 \\
1024-P-S10e-T & high & 54 & 81,36 & 88,89 & 84,96 \\
\end{longtable}\endgroup


\subsection{Truy vết nguồn và cách tổng hợp}
CSV tổng hợp được lưu ở \RepoPath{report/generated/phone-benchmarks/summary.csv}; các file \path{predictions.csv}, \path{class_metrics.csv} và \path{coverage.csv} nằm trong thư mục con \path{evaluate_model_phone}. Bốn CSV từng điện thoại nằm trong \path{evaluate_phone}; các bảng so sánh model và runtime ở \path{evaluate_baseline_pruning} và \path{evaluate_onnx_pte}. Mỗi dòng summary giữ nguồn JSON, hash nguồn, hash model/split, run ID và timestamp UTC để đối chiếu. Dự đoán từng ảnh và ma trận nhầm lẫn gốc vẫn được giữ trong CSV/JSON, không in hàng nghìn dòng vào báo cáo.

Bảy bảng được sinh bởi \RepoPath{report/scripts/build_phone_report_tables.py}; công cụ đối chiếu 32 JSON với hash nguồn, dự đoán, chỉ số, thời gian và support từng lớp trước khi xuất bảng. Bảng tóm tắt bốn phương án ở Chương~4 lấy trung bình chỉ số của tám phiên máy--runtime; số lỗi nặng cũng lấy trung bình nên có thể không nguyên. Các bảng trong phụ lục này giữ chỉ số riêng của từng phiên; không thay thế bằng một ma trận gộp các ảnh được dùng lại.
